\section{Data, diagnostic and analysis plan}
\label{sec:diagnostics}

\subsection{Data and apparent reproduction number}
\label{sec:diag-data}

We used the National Healthcare Safety Network (NHSN) weekly hospital respiratory-data snapshot retrieved on 20 October 2026 and restricted influenza admissions to the 50 US states and the District of Columbia \citep{cdc2026nhsn, frutos2025evaluation}. Hospital participation varied during the series. Reporting was voluntary from May to October 2024 before it became mandatory again from 1 November 2024 \citep{cms2024reporting}, and the national Respiratory Illnesses Data Channel was not updated from 3 October to 13 November 2025 \citep{cdc2026datachannel}. The analysis describes this surveillance system and does not estimate an effect of either event. Two external sources were added for the analyses of Section~\ref{sec:diag-amend2}: state resident populations from the Census Bureau Vintage 2025 estimates \citep{census2026popest}, and the issue history of the same NHSN admission series from the Delphi Epidata archive \citep{reinhart2021open}, which records the first published and every revised value of each state-week count from November 2024 onwards.

The windows were the early-growth phases of four seasons: 1 October--26 November 2022, 7 October--23 December 2023, 9 November--28 December 2024 and 4 October--27 December 2025. Each is the contiguous national segment with weekly growth of at least 1.05 before the first turning point after the season's admissible start; the 2024 start follows the return to mandatory reporting. The rule is applied identically to every season, and inference conditions on the selected windows. Sensitivity analyses use growth thresholds of 1.02 and 1.10 and fixed calendar windows for the first three seasons (8 October--26 November 2022, 14 October--16 December 2023 and 9 November 2024--18 January 2025). Missing or suppressed counts were kept missing rather than set to zero.

With $C_t$ the admissions in week $t$, the apparent reproduction number of a state-week is the growth-to-reproduction transform:
\begin{equation}
R_{A,t}=[\max\{0,\,1+(\theta/7)\log\{(C_{t+1}+0.5)/(C_t+0.5)\}\}]^{\alpha},\qquad
\alpha=(g_\mu/g_\sigma)^2,\quad\theta=g_\sigma^2/g_\mu.
\label{eq:ra}
\end{equation}
The transform follows the renewal relation between growth rate and reproduction number for a gamma-distributed interval \citep{wallinga2007how}. It is not a true reproduction number or a posterior estimate from a renewal model \citep{cori2013new, eales2024differences}. A row requires observed counts in weeks $t$ and $t+1$ and $C_t\ge5$, which leaves 1,748 state-weeks from 51 jurisdictions and 42 calendar weeks. The rows contain 1,969 distinct state-week counts in 210 runs of consecutive weeks.

The primary interval has mean $g_\mu=3.6$ days and standard deviation $g_\sigma=1.6$ days, the household serial-interval estimate of \citet{cowling2009estimation}. That estimate is a Weibull fit to 14 household transmission pairs; Eq.~\eqref{eq:ra} uses the gamma distribution with the same mean and standard deviation. A serial interval is not a generation interval \citep{vink2014serial}. Sensitivity analyses set the mean to 3.2 and 4.0 days and the standard deviation to 1.2 and 2.0 days; the lower mean is close to the intrinsic generation-time estimate from US household data \citep{chan2025estimating}.

\subsection{The diagnostic and its three summaries}
\label{sec:diag-definition}

For the apparent statistic $X=R_A$ with pooled mean $\bar X$, the relative squared deviation and its affine working model are:
\begin{equation}
L_{st}=\frac{(X_{st}-\bar X)^2}{\bar X^2},
\qquad \E(L_{st}\mid C_{st})=a+\frac{b}{C_{st}}.
\label{eq:affine}
\end{equation}
This is the empirical counterpart of $V(m)/\Rbar^2$ in Eq.~\eqref{eq:var-average}, with a realised count in place of the expected scale and an apparent statistic in place of $\Rhat_t$. It has the form of the additive overdispersion model of funnel plots \citep{spiegelhalter2005funnel, spiegelhalter2012statistical}. Three summaries of the fit follow.

The intercept $a$ is the count-independent relative dispersion: all relative variation in $X$ that does not scale with count. When reporting probabilities do not vary from week to week ($\tau=0$), $a$ estimates the relative variance of transmission, $\tilde v_R=v_R/\Rbar^2$. When reporting shocks vary between weeks, $a$ estimates $\tilde v_R+\tilde v_H$, where $\tilde v_H$ is the relative variance that the shocks induce in the ratio statistic (Appendix~\ref{app:ratio-shocks}). A reporting level that is constant over consecutive weeks cancels in the ratio, as it does in renewal estimators \citep{cori2013new, gostic2020practical}. By Proposition~\ref{prop:report-noise}, single-period cross-sectional moments cannot separate the two components. The estimator also need not weight them equally: when the shocks are drawn independently for each state-week, the fitted intercept recovers only part of $\tilde v_H$, and Section~\ref{sec:res-bias} quantifies the shortfall. The intercept is therefore read throughout as reflecting transmission and reporting variation, not as an estimate of their sum. It is also not an error floor for the current value: with $\tau=0$ that floor is zero, and with $\tau>0$ it is $\tilde v_H$. For the process mean, $a$ approximates the relative error that more cases cannot remove, to the extent that the affine form holds (Section~\ref{sec:res-form}).

The counting slope $b$ is read against two references. A delta expansion of Eq.~\eqref{eq:ra} under independent Poisson counts in consecutive weeks gives the independent-Poisson reference
$\bP=\overline{\kappa(g)^2[1+(C_t+0.5)/(C_{t+1}+0.5)]}$, a mean over rows, with $\kappa(g)=(g_\mu/7)/[1+\theta\log(g)/7]$ and $g=(C_{t+1}+0.5)/(C_t+0.5)$. Each row is evaluated at its own growth factor and normalised by $\bar X^2$, which approximates the expansion around $\bar X$ when the row's $R_A$ is close to $\bar X$; $\bP$ is therefore a plug-in normalisation, not an exact expectation. For the primary sample $\bP=0.4455$. Two alternatives show how much the normalisation matters. Replacing the unit normalisation by each row's ratio $X_{st}/\bar X$ gives 0.4199. The slope fitted to panels of independent Poisson counts drawn around the observed means, with the same transform, row filter and fit, gives 0.3977 (Table~\ref{tab:r2-poisson}). The term 1 is the Poisson variance of the denominator count. In a branching chain the denominator week contains the parents of the numerator week, so conditioning on the realised parents leaves only the binomial reporting part of that variance, $(1-\rho)/C_t$. With Poisson offspring and reporting probability $\rho$ this gives the branching-Poisson reference and its ratio to $\bP$:
\begin{equation}
b_{\mathrm{P,br}}(\rho)=\overline{\kappa(g)^2[(1-\rho)+(C_t+0.5)/(C_{t+1}+0.5)]},\qquad
r(\rho)=\frac{b_{\mathrm{P,br}}(\rho)}{\bP}\approx\frac{1+(1-\rho)g}{1+g}=1-\frac{\rho g}{1+g}.
\label{eq:branching-reference}
\end{equation}
Averaged over the observed rows, $r(\rho)$ is 0.944, 0.861, 0.722 and 0.443 for $\rho=0.1$, 0.25, 0.5 and 1 (Table~\ref{tab:ma_baselines}). The observed-denominator model of Section~\ref{sec:model-observed} has the structure of the branching-Poisson reference, and the known-scale estimator of Section~\ref{sec:model-branching} is its limit without denominator noise. The ratio $b/\bP$ asks whether counting variation exceeds independent-Poisson noise; $b/b_{\mathrm{P,br}}(\rho)$ asks the same for a branching chain, at a reporting probability that the data do not identify.

The third summary is the crossover count $b/a$, at which the fitted counting term equals the count-independent term. It is the empirical analogue of $\mx$ in admissions per week; $\mx$ refers to the expected count of a branching experiment. At $C=b/a$ a shrinkage estimator of the Fay--Herriot type places half its weight on the direct estimate (Eq.~\eqref{eq:shrink-weight}). Table~\ref{tab:rosetta} lists the conditions under which each summary identifies its theoretical counterpart and the direction of distortion when they fail.

\begin{table}[!tbp]
\centering
\scriptsize
\setlength{\tabcolsep}{3pt}
\caption{Mapping from theoretical inferential targets to empirical surveillance diagnostics}
\label{tab:rosetta}
\begin{tabularx}{\linewidth}{@{}p{0.20\linewidth}p{0.18\linewidth}p{0.28\linewidth}p{0.28\linewidth}@{}}
\toprule
Theoretical Target (Section~\ref{sec:model}) & Empirical Diagnostic (Section~\ref{sec:diagnostics}) & Exact Equivalence Conditions & Distortion Mechanism and Bias Direction \\
\midrule
Relative transmission variance: $\tilde v_R = v_R / \Rbar^2$ & Dispersion Intercept: $a$ & Constant reporting probability ($\tau = 0$); scale-independent transmission ($\gamma = 0$). & If reporting varies ($\tau > 0$), $a \approx \tilde v_R + \tilde v_H$. The intercept is an upper bound on $\tilde v_R$, not an error floor. Under scale dependence ($\gamma > 0$), $a < \tilde v_R$. \\
\addlinespace[2pt]
Counting error scale: $\tilde A = A / \Rbar^2$ & Counting Slope: $b$ & Independent Poisson limit ($\rho \to 0$ gives $b \approx \bP$); branching Poisson ($b \approx b_{\mathrm{P,br}}(\rho)$). & Inflated by offspring overdispersion ($k < \infty$, $b > b_{\mathrm{P,br}}$) or unmodelled scale-dependent heterogeneity ($\gamma > 0$). \\
\addlinespace[2pt]
Crossover scale: $\mx = A / v_R$ & Crossover Count: $b/a$ & Exact match between $A$ and $b$, and between $v_R$ and $a$. & Under reporting shocks ($\tau > 0$), $b/a$ is attenuated below $\mx$; under scale dependence ($\gamma > 0$), $b/a$ is inflated. $b/a$ is an in-sample balance point. \\
\bottomrule
\end{tabularx}
\par\smallskip
\begin{minipage}{\linewidth}\footnotesize\raggedright Formal correspondence between the theoretical parameters derived under the branching model in Section~\ref{sec:model} and the empirical summaries estimated by the cross-sectional dispersion diagnostic in Section~\ref{sec:diagnostics}. The table makes explicit the conditions under which empirical diagnostics identify theoretical quantities and the sign of the distortion when those conditions fail.\end{minipage}
\end{table}


\subsection{Estimation and inference}
\label{sec:diag-estimation}

The affine mean was fitted by iterated feasible generalised least squares with weights inversely proportional to the squared fitted mean, which minimises a gamma quasi-deviance; this is an estimating criterion, not a distributional claim. The primary interval is the two-way cluster-robust sandwich of \citet{cameron2011robust} for the stacked estimating equations of $\bar X$ and the affine fit, clustered by state and calendar week, with small-sample factors $G/(G-1)$ and critical values from the $t$ distribution on 41 degrees of freedom, one fewer than the number of weeks. The interval for $b/a$ is formed on the log scale. Two resampling checks accompany it: a calendar-week block bootstrap (1,999 draws) that resamples whole weeks with all their state rows, and a wild cluster bootstrap-$t$ over calendar weeks with six-point weights and the null $a=0$ imposed (9,999 draws) \citep{cameron2008bootstrap, webb2023reworking}. In the wild bootstrap, losses at or below $10^{-7}$ in a bootstrap sample are set to $10^{-7}$ before the restricted fit; this affected 13.5\% of bootstrap rows, and Table~\ref{tab:r2-clamp} reports the test under other thresholds and without the clamp. The state-by-week pigeonhole bootstrap \citep{owen2007pigeonhole} used for the supplementary analyses is conservative on this design: under the fresh-draw generators of Section~\ref{sec:diag-sim} its standard deviation of $\hat a$ was 1.55--1.78 times the Monte Carlo standard deviation (Table~\ref{tab:r2-bootstrap-sd}). Every interval quoted in the text is listed with its type and source in Table~\ref{tab:r2-interval-types}; all are 95\% intervals, and an interval without a stated type has the type of the nearest preceding interval of the same paragraph.

\subsection{Analysis plan and primary test}
\label{sec:diag-plan}

The analysis plan was written after an exploratory analysis of the same data, and it lists the estimates and the season-test $p$ values that were then known, so that each choice can be judged against them. It was fixed, with a checksum, before any analysis in Section~\ref{sec:results} was run. Because the data had been seen, the plan does not make the primary test confirmatory in the sense of an untouched hypothesis; what it fixes is the robustness rule for the intercept, the primary calibration of the season test and the obligation to report every specified analysis. The data vintage, the sample and $R_A$ with the primary interval were frozen.

The primary test (E1) is of $H_0\colon a=0$ against a two-sided alternative at $\alpha=0.05$, tested with the two-way sandwich $t$ statistic on 41 degrees of freedom. The claim that count-independent dispersion remains is called robust only if the primary test rejects and the 95\% lower limit of $a$ stays above zero in each of nine sensitivity analyses. Two vary the generation-interval mean (S-a1) and standard deviation (S-a2). Three change the row filters: count floors of 10, 20 and 50 (S-a3), the fixed calendar windows (S-a4) and growth thresholds (S-a5). Three remove data: transitions that touch the 2025 data-channel interruption (S-a6), one season at a time (S-a7) and rows with hospital reporting coverage below 90\% in either week (S-a8). The last (S-a9) removes one calendar week at a time to measure influence.

A within-week intercept (E3) is a different estimand, not a sensitivity analysis. It replaces $\bar X$ by the mean of the same calendar week computed without the row's own state, which removes common national variation by week; the week mean including the own row and a version that also rescales by the week mean are reported as alternatives. Its interval is the calendar-week block bootstrap.

The secondary test asks whether season intercepts differ. In a model with four season intercepts $a_s$ and a common slope, $H_0\colon\log a_1=\dots=\log a_4$ is tested with the Wald statistic $W_{\mathrm{CGM}}$ that uses the two-way sandwich covariance. The season label is constant within a calendar week, so it is assigned at the level of 42 weeks in four groups. With few clusters at the level of assignment, wild cluster bootstrap tests can be badly mis-sized and randomisation inference at that level is recommended \citep{mackinnon2018wild, mackinnon2023cluster}. The primary reference distribution is therefore a permutation of season labels across calendar weeks that keeps the number of weeks per season (9,999 permutations); it assumes exchangeable weeks and ignores serial dependence between adjacent weeks. The season difference is called significant only if this $p<0.05$, whatever the other calibrations show. All other calibrations are reported, and Holm-adjusted $p$ values over E1 and the season test are given. Season intercepts are interpreted only as reflecting transmission and reporting variation together (Proposition~\ref{prop:report-noise}).

The counting slope $b/\bP$, the crossover count and the predicted relative squared error $a+b/C$ at $C=10$, 37, 100 and 1,000 were specified for estimation, without tests.

\subsection{Simulation checks}
\label{sec:diag-sim}

The conversion of $a$ into a transmission standard deviation needs a mechanism. For a simulated mechanism with recovery ratio $r=\E(\hat a)/\tilde v_R$, the implied relative standard deviation of the reproduction number is $(a/r)^{1/2}$. Its interval combines the two-way interval of $a$ with the Monte Carlo 2.5--97.5\% range of $r$. The plan required that this quantity be reported as a grid over mechanisms and never as a single estimate, and that $a^{1/2}$ be described as a relative standard deviation of transmission and reporting variation combined, not of transmission alone. A mechanism reproduces the counting slope when its Monte Carlo 2.5--97.5\% range of $b/\bP$ contains the observed value. Synthetic panels are generated on the exact design of the data. The expected count level is calibrated to the observed median retained count of 45, and the heterogeneity amplitude is calibrated so that the mean fitted intercept equals the observed $a$ (Appendix~\ref{app:audit}). Five checks with fixed criteria were specified:
\begin{itemize}
\item SC1, two generations per week: a daily renewal process with the primary generation interval, Poisson or negative binomial ($k=0.3$) offspring and weekly binomial reporting at $\rho\in\{0.003,0.01,0.1\}$. The one-generation recovery map carries over if $\E(\hat a)/v_R$ lies in $[0.8,1.25]$ in every cell.
\item SC2 and SC3, reporting and offspring dispersion: a grid of $\rho\in\{0.001,0.003,0.01,0.03,0.1,0.25,0.5,1\}$ and $k\in\{0.1,0.3,1,3,\infty\}$ with independent transmission heterogeneity, 100 panels per cell. SC2 requires intercept recovery in $[0.9,1.1]$ in every cell; SC3 requires the mean of $b/\bP$ to lie in $[0.9,1.1]$ in every cell with $\rho\le0.01$. The range $\rho=0.001$--0.01 for the probability that an infection becomes a reported admission is assumed from burden estimates \citep{reed2015estimating, hayward2014comparative}, not estimated.
\item SC4 and SC5, interval validity: these use generators that redraw the heterogeneity field in every replicate. FD-S1 has independent heterogeneity, and FD-S8 has scale-dependent heterogeneity with state-week reporting shocks. FD-S9 adds to FD-S8 a national calendar-week share of 0.25, 0.5 or 0.75 of the heterogeneity; this family was defined after the data had been seen. SC4 requires coverage of $\E(\hat a)$ by the two-way interval of at least 0.90 under every generator. SC5 calls the gap between the observed standard error of $a$ and the generators' Monte Carlo standard deviation explained if their ratio lies in $[0.67,1.5]$ for at least one generator.
\end{itemize}

\subsection{Amendments and additional analyses}
\label{sec:diag-amend2}

The plan was amended twice, and both amendments are distributed with it. Amendment 1 recorded implementation details that the plan had left open: seeds, the restricted fit and loss clamp of the wild bootstrap, and the construction of the daily-renewal generator. It stated that these details were fixed before the full runs, but the run logs show that the full runs of the primary test, the season test and the simulation checks had finished before it was written. The details were already in the code used for those runs, and no evidence was found that any of them was changed after a result was seen. The season script was edited once after its first full run. The text of that edit could not be recovered, but every quantity logged by both runs is identical, including the primary permutation $p$ value, the other calibrations, the data sensitivities and the power analysis.

Amendment 2 was fixed, with a checksum, before any of the following analyses was run. Each is secondary or descriptive, none can change the verdict of the primary test, the significance level, the calibration or the data, and each was specified with its reading rule.
\begin{itemize}
\item Functional form. The observed mean loss in each decile of $C_t$, with a calendar-week block-bootstrap interval, is compared with the affine fit, a power law $cC^{-\gamma}$ and a three-parameter model $a+bC^{-\gamma}$. The intercept is said to survive a free exponent if the block-bootstrap interval of $a$ in the three-parameter model excludes zero; $\gamma=1$ is called compatible with the data if its interval contains one.
\item Model-free floor. Among rows with $C_t$ at or above the 75th, 90th and 95th percentiles, the mean loss minus the mean counting term, with either $\hat b$ or $\bP$ as the slope, estimates the count-independent component without extrapolation.
\item Population size. The intercept is made linear in the standardised log population $z$ of the jurisdiction, $a_0+a_1z$, with a common slope, and $a_1=0$ is tested with the two-way sandwich on 41 degrees of freedom. If the interval of $a_1$ lies below zero, statements about what more cases buy are restricted to jurisdictions of comparable size and the error at a given count is reported by population tertile. A log-link intercept, a slope that also varies with $z$, and separate fits by tertile are secondary.
\item Bias of the intercept. Under the fresh-draw generators with reporting shocks, and in a sweep of the transmission share of their variance (1, 0.75, 0.25), the relative bias $\E(\hat a)/(\tilde v_R+\tilde v_H)-1$ and the weight $w_H$ that solves $\E(\hat a)=\tilde v_R+w_H\tilde v_H$ are estimated. If every bias interval excludes zero, $\hat a$ is described as reflecting, not estimating, the total.
\item Generator realism. Six cells of a state-realistic generator anchor transmission to the smoothed observed epidemic curve, scale its variance by state and allow clustered admissions. Four tests compare the Monte Carlo range of each cell with the data: the slope ratio $b/\bP$ (T1), the 90th percentile of $C_t$ (T2), the ratio of pooled to within-state slope (T3) and the between-state coefficient of variation of the intercepts (T4). The generators reproduce the data if all four pass in one cell. Their design used the observed slope and between-state dispersion, and no inferential result rests on them.
\item Reporting data. The fit is repeated with counts scaled by the share of hospitals reporting in each state-week, and on the first-issue and latest-issue counts of the archive for the 661 rows whose first issue is provisional.
\item Calibration details. The wild bootstrap is repeated under four loss clamps and without one; the standard deviation of the pigeonhole bootstrap is compared with the Monte Carlo standard deviation of $\hat a$ (it is called conservative only if the ratio exceeds one under every fresh-draw generator); and $b/\bP$ is reported under the three Poisson references.
\end{itemize}
